\section{Tätigkeit}

\subsection{Erstellen des Netzproviders}

%TODO: #edit Ubuntu und PPPoE klingen eher wie eine Anleitung, als wie eine "Dokumentation"
	\subsubsection{Ubuntu}
	Wie in \Verweis{subsec:ubuntu} beschrieben wird Version 14.04 LTS
	der Linuxdistribution Ubuntu eingesetzt. Das bei der Herstellerwebsite\footnote{\url{http://releases.ubuntu.com/14.04/}},
	für Serverbetrieb erhältliche Image wird in einer virtuellen Maschine betrieben.
	Diese weist folgende Eigenschaften auf:

	\begin{table}[h]
	\begin{tabularx}{\textwidth}{|l|X|}
    	\hline
      		\textbf{Gerät} & \textbf{Typ} \\
    	\hline
    	\hline
    		CPU & 1x 64 Bit \\
    	\hline
    	\hline
    		RAM & 1x 1024 MB \\
    	\hline
    	\hline
    		HDD & 1x 10 GB \\
    	\hline
    	\hline
    		NIC & 2x Bridged (1x auf enp0s25; 1x auf wlp3s0) \\
    	\hline
    	\end{tabularx}
       	\caption{Eigenschaften der virtuellen Maschine}
       	\end{table}

    	Aufgrund der Tatsache dass nur ein einziger Client über den Zugang von \textit{PPPoE}
    	zu erwarten ist, sind die Anforderungen an CPU, Arbeits- und Festplattenspeicher
    	recht gering. Wichtig ist das Vorhandensein zweier Netzwerkschnittstellen, die
    	auf jeweils eine physikalisch Schnittstelle geleitet werden. Dies ist notwendig,
    	da sowohl die \ac{MS} über eine Ethernetverbindung mit der virtuellen Maschine angebunden sein wird,
    	als auch eine Dauerhafte Internetverbindung mit der verbleibenden Schnittstelle
    	aufrecht erhalten werden muss. So ist es möglich nach erfolgreicher Authentifizierung
    	des \ac{MS} den Zugang zum Internet freizuschalten.

    	\subsubsection{PPPoE}
    	Die Einrichtung von \textit{PPPoE} basierend auf der vollständigen Ubuntuinstallation realisiert werden. Hierzu wird die aktuelle
    	Version 3.11 auf der Herstellerwebsite\footnote{\url{https://www.roaringpenguin.com/products/pppoe}}
    	bezogen. Zusätzlich wird das Ubuntu-Paket zur Unterstützung
    	des Point-To-Point-Protocol (als \textit{ppp} erhältlich) benötigt. Sind diese beiden Grundvoraussetzungen
    	abgedeckt wird der Roaring Penguin PPPoE-Server ohne spezielle Änderungen aus dem Quellcode kompiliert.

    	In der Datei \textit{/etc/ppp/pppoe-server-options} werden die notwendigen Parameter gesetzt. Anzugeben
    	sind die DNS-Server sowie die CHAP-Authentifizierung. Der User für die Authentifizierung
    	wird unter \textit{/etc/ppp/chap-secrets} definiert.

    	Der zu verwendende Adresspool an IP-Adressen ist ebenfalls (unter \textit{/etc/ppp/allip}) anzugeben.

    	Abschließend fehlt noch die Deklaration der zu verwendenden Netzwerkschnittstelle (unter \textit{/etc/network/interfaces}),
    	die zum Betrieb mit PPPoE eingesetzt werden soll. Ihr muss eine valide IP-Adresse aus dem definierten
    	Pool sowie eine Subnetzmaske zugewiesen werden.

    	Alle Konfigurationsdateien befinden sich im Anhang. %TODO: #edit Anhang pflegen

    	Damit die Internetverbindung nach erfolgreicher Authentifizierung des \ac{MS} auch
    	korrekt von Provider freigeschaltet und weitergeleitet wird, muss eine IPTables-Regel
    	folgernder Form eingerichtet werden:

    	\begin{lstlisting}[language=bash]
		$ iptables -t nat -A POSTROUTING -s 192.168.178.0/24 -o enp0s25 -j MASQUERADE
	\end{lstlisting}

	Sie nimmt den Verkehr auf dem Device \textit{enp0s25} an und maskiert (vgl. \ac{NAT}) diesen
	für den weiteren Betrieb (wie es bei einem realen Provider auch wäre). Dies geschieht
	für eingehenden sowie ausgehenden Verkehr. Die auf den zweiten Netzwerkadapter geleitete
	Netzwerkschnittstelle bleibt mit Defaultwerten konfiguriert und fungiert als Zugangspunkt
	zum Internet.

	Der Server kann auf der gewünschten Netzwerkschnittstelle mit einer eigenen IP-Adresse
	gestartet und genutzt werden:
	\begin{lstlisting}[language=bash]
		pppoe-server -C isp -L 192.168.178.254 -p /etc/ppp/allip -I enp0s25
	\end{lstlisting}

        Korrekt konfigurierte Clients, die über die Zugangsdaten verfügen können sich über
        diese Verbindung anmelden und Internetzugriff erhalten.

	\subsubsection{Implementierung AES}
	\label{implementierung-aes}
	Die Blockchiffre AES wurde in der Datei \emph{rijndael.c} geschrieben. Der Code wurde in drei
	logische Teile unterteilt, nämlich die Variablen\-deklaration, die Main-Funktion, die die
	einzelnen Runden steuert und Funktionen, die die konkrete Implementierung der jeweiligen
	Transformation bereitstellen. Des weiteren wurde ein eigener Variablentyp \emph{u8} definiert,
	der eine verkürzte Form des unsigned integer ist. \\
%%TODO: #anhang C-Code hinzufügen
	Der gesamte Code steht im Anhang ??? und wird in den folgenden Paragraphen näher erläutert.

		\paragraph{Variablen}
		Die zu Beginn initialisierten Variablen sind einige Zählervariablen, die für spätere Schleifen
		reserviert werden. Zusätzlich werden jedoch auch einige Konstanten definiert. Da keine
		Optimierung bezüglich des Speicherplatzverbrauchs getätigt werden muss, wird die Substitutions-
		Tabelle, also die S-box, nicht zur Laufzeit erst berechnet, sondern zu Beginn definiert. Ähnlich
		sind die Rundenkonstanten für die Generierung des Rundenschlüssels schon vorher in einer Konstante
		definiert. Des weiteren existiert die \emph{mixMatrix}, genauer gesagt ein 4 x 4-Array, welches die
		MixColumns-Transformation vereinfachen wird. \\
		\emph{tmp} ist ein Hilfs-Array, in \emph{key} wird der Chiffrierschlüssel gespeichert, in \emph{roundKey}
		die Rundenschlüssel für alle zehn Runden, also inklusive der 1. Runde, und in \emph{state} wird der Zustand
		hinterlegt. 
 
		\paragraph{encrypt}
		Die Funktion, die vom Milenage-Algorithmus aufgerufen wird ist \emph{encrypt} und
		braucht drei Eingabeparameter. \emph{input} ist der Textblock, \emph{keyStr} der
		Schlüssel und output die Variable, in die der chiffrierte Status am Ende geschrieben
		werden soll. \\
		Da Textblock und Schlüssel zu erst noch in ein Array konvertiert werden müssen
		wird bei beiden erst convert2array aufgerufen. Wie die Funktion erkennen lässt
		konvertiert sie den Inhalt des ersten Eingabeparameter in ein Array und speichert das
		Ergebnis in der zweiten übergebenen Variable. Entgegen der Spezifikation wird das
		der Textblock und der Schlüssel aber nicht zeilenweise, sondern spaltenweise in das
		Array umgewandelt.

		\begin{lstlisting}[language=C]
		void convert2array(u8 input[16], u8 output[4][4]) {
    			for (j = 0; j < 4; j++) {
        			for (i = 0; i < 4; i++) {
            				output[i][j] = input[(j * 4 + i)];
        			}
    			}
		} 
        	\end{lstlisting}

		\emph{j} ist die Spalte in dem Array und \emph{i} die Zeile. Durch die einfache
		Rechnung $j*4+i$ kann das Array nacheinander durchgegangenen werden und spaltenweise
		abgespeichert werden. 
		%Die Art der Speicherung hat das Programmieren der shiftRows-
		%Transformation vereinfacht, und andere Transformationen nicht erschwert. Die
		%Umwandlung des Textblocks und Schlüssel in ein Array und zurück wäre jedoch einfacher,
		%da jediglich die Arraystruktur verändert werden müsste, aber nicht die Speicher\-
		%struktur, wie in dem hier umgesetzten Fall.

		Anschließend wird gemäß der Spezifikation der Rundenschlüssel generiert, allerdings
		wie schon erwähnt für alle zehn Runden und nicht für jede Runde einzeln. Wie im
		\Anhang{abb:funktion_aes} gezeigt, wird dann der erste Rundenschlüssel mit dem
		Zustand verknüpft. \\
		Im Anschluss folgen die ersten neun Runden durch eine einfache for-Schleife. Die
		folgenden beiden verschachtelten Schleifen sind schon die SubByte-Transformation,
		welche nicht in eine eigene Funktion ausgegliedert ist. \emph{getSboxValue} dient
		nur zur besseren Lesbarkeit des Codes. \\
		Nachfolgend werden auch die anderen Transformationen durchgeführt und nach den neun
		Runden die zehnte ohne eine MixColumns-Transformation.

		Der finale Zustand wird dann wieder in eine Zeichenkette konvertiert und
		zurückgegeben.

		\paragraph{generateRoundKey}
		In der ersten doppelten for-Schleife von generateRoundKey wird der Chiffrierschlüssel
		in den ersten Block des roundKey übertragen. \\
		Anschließend werden die restlichen Rundenschlüssel gemäß Spezifikation generiert.
		Um die spaltenweise Rotation durchzuführen, wird das erste Byte der Spalte, das von
		oben nach unten rotiert, in das temporäre Array geschrieben. Anschließend
		überschreibt der zweite Wert den ersten, dann der dritte Wert den zweiten und im
		vorletzten Schritt das vierte Byte in der Spalte. Da nun der Wert im ersten Byte
		nicht mehr der selbe wie vor der Rotation ist, wird das vierte Byte nun mit dem
		Wert aus dem temporären Speicher gefüllt.

		Da die erste Spalte des Arrays anders berechnet wird als die folgenden, gibt es
		zwei Schleifenblöcke. Der erste berechnet die erste Spalte und der zweite Block,
		die restlichen drei, da diese nach dem gleichen Schema berechnet werden.

		\paragraph{shiftRow}
		Die Funktion shiftRow ist etwas komplexer, auch wenn sie nur wenige Zeilen lang ist.
		Die erste Schleife zählt die Spalten des Zustands (\emph{state}) durch und die
		beiden folgenden Schleifen die Zeilen. \\
		In der ersten inneren Schleife wird die aktuelle Zeile komplett in ein temporäres
		Array gespeichert und in der zweiten Schleife findet die Rotation statt.

		\begin{lstlisting}[language=C]
		for (i = 0; i < 4; i++) {
			for (j = 0; j < 4; j++) {
		        	tmp[j] = state[i][j];
		        }
	        	for (j = 0; j < 4; j++) {
	            		state[i][j] = tmp[(j + i) % 4];
	        	}
    		}
		\end{lstlisting}

		In der zweiten inneren for-Schleife wird dem Zustands-Array Spalte für Spalte und
		Zeile für Zeile ein neuer Wert zugewiesen, der aus dem temporären Array geliefert
		wird. Durch die Addition von \emph{j} und \emph{i} wird bei der ersten Zeile auch
		mit dem ersten Byte angefangen, bei der zweiten Zeile mit dem zweiten Byte und so
		weiter. Das Byte mit dem angefangen wird, wird aber immer an die erste Stelle des
		Zustands-Array gespeichert. Damit wird die Rotation um null, eins, zwei und drei
		Bytes ermöglicht. Da aber zum Beispiel in der dritten Zeile das
		das erste Byte rechts wieder reingeschoben werden soll, ist die modulo 4-Operation
		nötig. \\
		$i \% 4$ reicht übrigens nicht aus, da dann zwar beispielsweise beim dritten Byte
		in der dritten Zeile begonnen werden würde, aber nicht die durch die
		anderen Bytes in der Zeile iteriert wird. 

		\paragraph{mixColumn}
		In mixColumn findet eine umgekehrte Herangehensweise statt als in shiftRow.
		mixColumn iteriert erst durch die Zeilen und dann duch die Spalten. Dementsprechend
		wird in das temporäre Array die aktuelle Spalte geschrieben. \\
		Anschließend wird gemäß der Rijndael-Definition die MixColumns-Tranformation
		durchgeführt für jedes Byte in jeder Spalte. Das Array \emph{mixMatrix} gibt vor,
		ob die Funktion $T_2$ oder $T_3$ angewendet werden muss. In der Implementierung
		wurde jedoch nur eine t2-Funktion geschrieben und $T_3$ wird mit Hilfe der 
		t2-Funktion umesetzt, ähnlich der mathematischen Definition in Kapitel
		\Verweis{milenage-funktion}.

	\subsubsection{Implementierung Milenage}
	\label{implementierung-milenage}
	Die Implementierung des Milenage Algorithmus wurde nicht auf Basis der Funktionen
	strukturiert sondern der Ausgabeblöcke. Die Funktionen \emph{f2} und \emph{f5} sind
	deshalb in einer Funktion zusammengefasst, ebenso wie \emph{f1} und \emph{f1*}. Die
	beiden übrigen Funktionen \emph{f3} und \emph{f4} wurden hingegen gar nicht implementiert,
	da diese nicht für den reinen Authentifizierungsvorgang nötig sind und deshalb nicht
	gefordert sind. \\
	In den folgenden Paragraphen werden wieder die Hauptfunktionen erklärt für ein besseres
	Verständnis des Codes.
	
		\paragraph{Variablen}
		Ähnlich wie im AES-Programmcode sind auch bei Milenage einige Variablen zu Beginn
		deklariert. Dazu gehören alle Werte, die SIM-Karrten spezifisch sind und in der
		Realität in einer Datenbank gespeichert werden für jeden Netzteilnehmer. So sind also
		die Rundenkonstanten c1 bis c5, SQN, AMF und OPc schon fest definiert. Zusätzlich sind
		die Adressräume für die Werte, die generiert werden reserviert, wie AK oder RES und \emph{i}
		als Zählervariable.
		
		\paragraph{f1}
		Die Funktion benötigt als Eingabeparameter den Schlüssel K und den RAND. Der RAND wird in der
		Funktion genRand generiert auf Basis der aktuellen Zeit. \\
		Für Demozwecke werden dann zu Beginn der Funktion die beiden Werte SQN und AMF ausgegeben, so
		wie in der restlichen Funktion auch andere Parameter in der Konsole ausgegeben werden. \\
		Die beiden folgenden Schleifen sind zur Generierung des IN1, welcher ein konkatenierter Wert aus
		SQN und AMF ist. Um nicht in vier Schleifen abwechselnd durch beide Werte zu iterieren wird nur
		jeweils einmal durch beide Werte iteriert. Im Falle von SQN wird das erste Byte also direkt an den
		Anfang des IN1-Arrays und an achte Stelle geschrieben. Ähnlich funktioniert auch das hinzufügen von
		AMF zu IN1. \\
		Die Variable \emph{toEncrypt} ist nur eine temporäre Variable in der der OPc mit dem RAND verknüpft
		wird über XOR. Diese wird dann verschlüsselt mit der in Kapitel \Verweis{implementierung-aes} erklärten
		Funktion encrypt und der verschlüsselte Wert ist dann der TEMP-Wert. Dieser wird nachher auch von den
		anderen Funktionen benötigt, aber da \emph{f1} die erste Funktion ist die aufgerufen wird, wurde die
		Generierung der TEMP-Funktion nicht ausgegliedert in eine eigene Funktion.
		
		Anschließend fängt die Berechnung des ersten Blocks OUT1 an aus dem sich MAC-A und MAC-S ableiten. Nach
		der Verknüpfung von OPc und IN1 wird rotiert. Anders als bei AES und der shiftRows-Transformation ist die
		Rotation nicht byteweise, sondern bitweise. Die Strings sind aber als Byte-Arrays gespeichert, weshalb sie
		für eine bitweise Rotation erst zu einem binären Array konvertiert werden. Diese Konvertierung geschieht
		über eine Bitmaske. In jedem Byte des \emph{hexArr} werden alle acht Bit durchiteriert und je nach dem, ob
		da ne 1 oder 0 steht wird diese in das \emph{binArr} geschrieben. Die nachfolgende Rotation wird in einem
		der folgenden Paragrafen noch näher erläutert. Die Rückkonvertierung des binären Arrays in das
		Byte-Array erfolgt wieder über Bitoperationen. Das entsprechende Byte wird erst um 1 Stelle nach links
		verschoben und es wird rechts mit einer 0 aufgefüllt. Durch die Oder-Operation mit dem Bit-Array wird diese
		zu einer 1, wenn eine 1 im Bit-Array steht oder bleibt eine 0.
		
		Die restliche f1-Funktion verläuft wie in der Spezifikation definiert. Der rotierte OUT1 wird mit der
		Konstante c1 und TEMP verknüpft, verschlüsselt und zum Schluss noch mit OPc verknüpft. Die beiden letzten
		Schleifen geben nur noch MAC-A und MAC-S aus.		
		
		\paragraph{f2\_5}
		Die Funktion f2\_5 braucht zwei Eingabeparameter, nämlich den Schlüssel und \emph{response\_arr} welches
		die Werte enthält, die nach der Registrierungsanfrage an die SIM-Karte zurückgeschickt werden. \\
		Ansonsten ist die Funktion sehr ähnlich der schon beschrieben f1-Funktion aufgebaut. Jedoch wird IN1
		nicht benötigt und TEMP wurde schon generiert, weshalb diese Werte in f2\_5 nicht nochmal generiert
		werden müssen. So wird also gemäß der Spezifikation TEMP mit OPc verknüpft, dieser Wert rotiert und mit
		der Konstanten c2 verknüpft. Das Ergebnis ist OUT2, aus welchem sich AK und RES ableiten, welche in jeweils
		einer for-Schleife auch wieder in der Konsole ausgegeben werden. Da RES auch an die USIM geschickt wird als
		Antwort, wird diese in das \emph{response\_arr} geschrieben, wofür die Standardfunktion sprintf genutzt wird.
		
		\paragraph{f5star}
		Diese Funktion benötigt ebenfalls zwei Werte. Das eine ist wieder der Schlüssel und das andere ist der SQN
		verknüpft mit AK, der von der SIM-Karte geschickt wird für die Resynchronisation. Die Funktion berechnet erst
		den AK aus \emph{f5*} anhand der Definition des 3GPP. Dafür wird wieder OPc mit TEMP verknüpft über die
		XOR-Operation, das ganze rotiert, mit c5 verknüpft und verschlüsselt. Zum Schluss noch einmal wieder mit
		OPc verknüft und damit ist OUT5 fertig berechnet.
		
		Die erste Hälfte, die AK ergibt wird nun benötigt, um die SQN aus der Variable \emph{sqn\_ak} zu ermitteln.
		Diese SQN ist allerdings keine, die für die Berechnung der neuen Authentifizierungsvektoren genutzt werden
		kann, da sie schon benutzt wurde und der Netzprovider deshalb eine neue SQN generieren muss. Bei der USIM
		in diesem Projekt wird eine komplett zählerbasierte SQN verwendet mit einer Arraygröße von 32.
		Wie in Kapitel \Verweis{milenage-funktion} näher erläutert, bedeutet das, dass die hinteren fünf Bits der
		IND sind und die vorderen 43 die SEQ. Da die SQN ein Bytearray ist, wird aus dem letzten Byte mit einer
		Bitmaske die IND extrahiert und die vorderen drei Bits des letzten Bytes als SEQ abgespeichert. Bevor
		die SEQ hochgezählt wird, wird der IND erhöht. Da das SQN-Array nur 32 Werte speichert, wird IND nicht
		beliebig groß und fängt durch modulo 32 wieder bei 0 an, wenn die 31 überschritten wird. Das letzte Byte
		wird nun auf 0 gesetzt und die letzten fünf Bits gleich dem neuen Wert von IND gesetzt.

		Im zweiten Schritt wird auch SEQ um eins erhöht, aber da SEQ beim sechsten Bit anfängt wird SEQ nicht mit
		1 sondern mit $32_{10}$ ($00100000_2$) addiert. Wenn jedoch die vorderen drei Bits des letzten Bytes $111_2$ sind
		und jetzt mit 1 addiert wird, muss ein Übertrag in das vorletzte Byte vorgenommen werden. Diesen Überlauf
		überprüft die if-Abfrage und addiert im Falle eine 1 auf das vorletzte Byte. Unabhängig vom Überlauf wird
		der Wert von \emph{seq} auch wieder in das letzte Byte durch die oder-Operation übertragen.
		%%TODO: #edit überprüfen in wie vielen Anfragen eigentlich auch der Überlauf des 5 Bytes stattfinden kann und Code eventuell anpassen

		\paragraph{rotWord}
		%%TODO: Muss im Code in die milenage.c verschoben werden
		Diese Funktion nutzt einen effizienten Algorithmus um ein Array oder eine Zeichenkette zu rotieren. Die
		Idee dahinter ist, dass	wenn wie in \Abbildung{rotWord} eine Zeichenkette 10 Zeichen hat und um 3 Zeichen
		nach links rotiert werden soll, dann werden die ersten drei Zeichen invertiert, anschließend die letzten
		sieben Zeichen und in einem dritten Schritt der gesamte String.
		%%TODO: #edit richtige Abbildung einfügen
		\begin{figure}[htp]
			\begin{center}
				\fbox{\includegraphics[width=250pt]{rotWord}}
			\end{center}
			\caption{Schematische Darstellung des Rotations-Algorithmus}
			\label{fig:rotWord}
		\end{figure}
		
		In der konkreten Implementierung wird zum invertieren die Funktion reverse aufgerufen, die als
		Übergabeparameter ein Array übergeben bekommt und die Größe des Arrays, dass invertiert werden muss.
		Das vordere und hintere Element werden dann mit Hilfe der \emph{tmp}-Variable ausgetauscht, bis die Schleife
		in der Mitte des Arrays angekommen ist und das Array damit invertiert wurde.

	\subsection{Erstellen des UE}
		\subsubsection{Aufsetzen des Raspberry Pis}
        \label{subsubsec:installpi}
			\paragraph{Raspbian} wird über die von der Webpräsenz\footnote{\url{https://www.raspberrypi.org/downloads/raspbian/}}
            des Projekts gehosteten vorgefertigten Builds aufgesetzt.

            Die SD-Karte muss nicht weiter formatiert und für die Installation vorbereitet werden. Das
            bereitstehende ISO-Image ist ein Abbild eines bereits korrekt formatierten sowie installierten
            Raspbian (Jessie). Nach kopieren (blockweise) des Images auf die SD-Karte kann von dieser direkt
            gebootet werden.

            Nach dem initialen Bootvorgang können Userspezifische Änderungen durchgeführt werden. Hierzu
            gehört das anpassen von Netzwerk- und Hostinformationen. Um später eine Administation via
            \ac{ssh}-Remotezugang zu ermöglichen wird auch dieser in der entsprechenden Konfigurationsdatei
            zum automatischen starten aktiviert.

            Für den weiteren Betrieb müssen Pakete aus dem offiziellen Resporitory
            nachinstalliert werden. Spezielle Pakete, die dort nicht verfügbar sind werden
            manuell kompiliert.

            \paragraph{SIM-Kartenleser} Der benötigte (herstellerspezifische)
            Cherry-Treiber für den Kartenleser TC1100 ist nicht im offiziellen Repository enthalten, kann jedoch
            manuell kompiliert werden. Auf der Webpräsenz des Herstellers\footnote{\url{http://www.cherry.de/cid/download.php}}
            wird ein Treiber für den Betrieb mit 64-Bit Linuxdistributionen zu Verfügung gestellt.
            Selbiger kann nach Installation der nativen \textit{pcsc}-Pakete (aus dem Repository)
            kompiliert und eingesetzt werden.

            Sind diese Softwarepakete eingerichtet kann der Leser via USB verbunden, die SIM-Karte eingesetzt
            und erste Kommunikationsversuche durchgeführt werden. Hier bietet sich der vom
            Treiber mitgelieferte Scanmechanismus für neue Devices an. Da der Treiber
            Hotplugfähig ist kann direkt mit der SIM-Karte gearbeitet werden, sobald diese
            eingesetzt wird.

			\paragraph{pysim} kann von Kevin Prince's Githubrepository\footnote{\url{https://github.com/kevinprince/pysim}} bezogen werden. Es benötigt ebenso
            wie andere nachfolgend genannte Python-Skripte noch das Paket \textit{pysim}, welches
            in den offiziellen Raspbian Paketquellen zur Verfügung steht. Nach erfolgreicher Installation
            von Pysim und pullen des Skriptes wird selbiges über Python verwendet.

            \paragraph{osmo-sim-auth} ist vom Gitrepository\footnote{\url{}http://cgit.osmocom.org/cgit/osmo-sim-auth} des Projekts Osmocom zu  beziehen. Genau wie auch pysim kann dieses nach dem pullen
            über Python verwendet werden. Entweder als Bibliothek oder direkt im Interpreter.

	\subsection{Integration mit PPPoE}
        Um den Verbindungsaufbau zwischen \ac{MS} und \ac{AuC} korrekt abzuwickeln muss sowohl der Client
        als auch der Server für \ac{PPPoE} konfiguriert sowie gestartet sein.

		 \subsubsection{Server}
         Hier muss initial der Daemon für den \ac{PPPoE}-Server gestartet sein. Anzugeben ist das gewünschte
         Netz mit zugehöriger Netzwerkschnittstelle. Ebenso muss die IPTables-Regel zur Weiterleitung
         des Internetverkehrs gesetzt sein.

         \subsubsection{Client}
         Hier müssen einmalig über das Tool \textit{pppoe-conf} die Parameter und Authentifizierungsdaten
         für die \ac{PPPoE}-Verbindung eingepflegt werden.

         Das Tool bietet die Möglichkeit selbige Informationen
         zu speichern und automatisch den Verbindungsaufbau mit dem Server aufzunehmen. Bei Neustart
         des Systems oder Verbindungsabbruch (regulär oder durch technische Zwischenfälle) greift dieser
         Mechanismus ebenfalls.
         
\subsection{Implementierung der Server-Client Architektur}
\label{server-client-architektur}
	\subsubsection{Implementierung des Servers}
    	\label{subsec:impl_server}
	Das Hauptprogramm, das am Ende in main.c gestartet wird, ist der eigentliche Netzprovider beziehungsweise
	die Authentifizierungsstelle. Damit diese Daten dynamisch empfangen kann, wird ein Serversocket aufgebaut
	mit dem sich Clients, also die USIM, verbinden und kommunizieren können. Um diesen Socket aufzubauen werden
	die beiden Bibliotheken \emph{sys/socket} und \emph{arpa/inet} benötigt. Zusätzlich werden zu Beginn wieder
	einige Variablen deklariert, die benötigt werden, um eine gute Kommunikation mit als Server mit der SIM-Karte
	zu ermöglichen. \\
	Die Konstanten mit Präfix ANSI wurde definiert, damit die Konsolenausgaben farblich sind, aber die Lesbarkeit des
	Codes darunter nicht leidet.
	
		\paragraph{Kommunikationsaufbau}
		Beim Start des Programms muss dem Server ein Socket zur Verfügung gestellt werden, was mit der Funktion
		socket erfolgreich geschieht, sofern sie keinen Fehlercode zurückgibt, der in der if-Abfrage abgefangen wird.
		Anschließend muss	der Serversocket konfiguriert werden. Er wird mit der Funktion \emph{bind} so eingestellt, dass er
		auf jede IP-Adresse mit Port 12345 hört. Nachdem dies konfiguriert ist, wartet das Programm auf eine Verbindungsanfrage
		durch die Funktion \emph{listen}. Neben der Socket-Beschriebung muss der Funktion listen mitgeteilt werden, wie viele Verbindungen
		der Server maximal zulässt. Da in diesem Projekt nur mit einer SIM-Karte gearbeitet wird, ist der Wert also auf 1 gesetzt.
		Nach diesen Schritten akzeptiert der Server einen Client und speichert die Daten des Clients in \emph{client\_sock}.
		Wenn ein Client nun eine Verbindung aufbaut, wird dies vom Server registriert und die Infos des Clients abgespeichert,
		damit der Server weiß mit wem er kommuniziert. \\
		In der while-Schleife wird nun darauf gewartet, dass die USIM eine Nachricht schickt, die mit der Funktion recv empfangen
		wird. Die Funktion recv benötigt den Client von dem eine Nachricht erwartet wird, sowie ein Array in den die Nachricht dann
		geschrieben werden soll. Außerdem wird angegeben, wie lang die Nachricht sein wird. Diese Zahl darf maximal
		so groß sein, wie die Größe des Arrays, in dem die Nachricht gespeichert wird. Zusätzlich können noch Flags
		gesetzt werden, was in diesem Fall aber nicht nötig ist. \\
		Wenn die Funktion recv nun einen Wert zurückgibt kann dieser -1, 0 oder größer 0 sein, wobei letzteres die
		Länge der empfangenen Nachricht bedeutet. Die anderen beiden Fälle werden am Ende dieses Kapitels erklärt.
		
		\paragraph{Verarbeiten der Nachrichten}
		Wenn eine Nachricht empfangen wurde ist sie im übergeben Array \emph{client\_message} eingetragen. Dieses
		Array ist eine Zeichenkette und kein Array aus unsigned chars. Außerdem sind in dieser Nachricht einige
		Werte konkateniert gespeichert, weshalb diese nun getrennt werden. Zu erst wird das Statusbyte in eine
		temporäre Varaible gespeichert mit Hilfe der Standardfunktion strcpy und anschließend wird der Wert
		umgewandelt in eine Dezimalzahl mit strtol. Die Konsolenausgabe ist abhängig davon, ob eine erste Anfrage
		oder eine Resynchronisation nötig ist. \\
		Als nächstes wird die IMSI extrahiert, welche unbearbeitet in \emph{response\_message} später wieder
		eingetragen wird und der Status in \emph{response\_message} wird standardmäßig auf 50 gesetzt. 50 ist
		dabei die hexadezimale Darstellung des Zeichens ``2''.
		
		Wenn die Nachricht die empfangen wurde den Status 3 hat, dann bedeutet das, dass Resynchronisiert werden
		muss. In dem Falle wurde von der USIM der AUTS geschickt und dieser wird in entsprechende Variable übertragen.
		Nach der Extraktion wird mit AUTS die Funktion f5star aufgerufen, um den Resynchronisations AK zu berechnen
		sowie die neue und richtige SQN zu kalkulieren. Des weiteren wird der Status der Antwortnachricht auf 52, also ``3''
		gesetzt.
		
		Im letzten Schritt wird RAND generiert, \emph{f1}. \emph{f2} und \emph{f5} berechnet, sowie aus den generierten
		Werten der AUTN gebildet. Die Werte werden alle in der Antwortvariable eingetragen, damit dieser dann an die USIM
		geschickt werden kann. Das passiert über die Funktion write, die wie recv die Daten des Clients braucht, die eigentliche
		Antwortnachricht und die Länge dieser Nachricht. Für die bessere Visualisierung wird die Antwort am Ende auch in
		der Konsole ausgegeben.
		
		\paragraph{Beenden des Servers}
		Wie schon erwähnt kann recv, auch die Werte 0 und -1 zurückgeben. Im Falle von 0 bedeutet das, dass der Client
		die Verbindung unterbrochen hat und deshalb keine Kommunikation mehr möglich mit diesem. Bei -1 hingegen bedeutet
		es, dass die Nachricht des Clients nicht gelesen werden konnte. \\
		In beiden Fällen wird das Programm und damit der Server beendet.

    \subsubsection{Implementierung des Clients}
    Wie bereits erläutert existiert neben dem Server (vgl. \Verweis{subsec:impl_server}) auch ein Client,
    damit die Kommunikation wie erwünscht realisiert werden kann. Hier handelt es sich um die Komponente,
    die Zugriff auf die SIM-Karte hat (\emph{UE}) und alle notwendigen Informationen an den Server sendet.
    Alle notwendigen Funktionen sind entweder auf der SIM-Karte persistent gespeichert oder werden durch
    das Triggern der integrierten Funktionen generiert.

    Der Client wird in Python entwickelt und verwendet grundlegend die Bibliothek \textit{osmo-sim-auth}
    (vgl. \Verweis{subsec:osmosim}). Darüber hinaus kommen weitere  Bibliotheken zum Einsatz, die in
    Python den Kommunikationsaufbau mit Chipkarten ermöglicht (card-Utils, card-Sim und card-USIM).

        \paragraph{Aufruf der Authentifizierung} Diese Funktion wird im Clientprogramm im Abschnitt
        \textit{handle\_usim} definiert. Es werden externe Informationen (falls vorhanden) für den Aufruf
        gesammelt und entsprechend an die SIM-Karte gesendet. In jedem Fall wird ein Authentifizierungsvorgang
        für eine USIM-Karte initiiert. \\
        Abhängig von gegebenen Information ist die SIM-Karte daraufhin in der Lage Ergebnisvektoren
        zu berechnen und zu liefern. Hier gibt es drei verschiedene Möglichkeiten:
            \begin{addmargin}[1em]{2em}
            Die Form der Anfrage war inkorrekt. D.h. sowohl der Wert für AUTN, als auch der Wert für RAND entsprechen
            den standardisierten Vorgaben. Es ist also notwendig die Anfrage mit korrekten Werten erneut
            durchzuführen.

            Die Form der Anfrage war korrekt, jedoch ist eine Resynchronisierung der übermittelten SQN erforderlich.
            In diesem Fall muss der Aufruf des Authentifizierungsvorgangs nach erfolgreicher neuberechnung der SQN
            erneut durchgeführt werden.

            Die Form der Anfrage war korrekt, sowie die Beschaffenheit der übermittelten SQN. Ist dies der Fall,
            kann die SIM-Karte ohne weiteres die Vektoren Ck, Ik und RES generieren. Diese werden entsprechend
            ausgegeben.
            \end{addmargin}
        Der Aufruf des Authentifizierungsvorgangs kann im Abschnitt \textit{main} über \textit{authenticate}
        eingeleitet werden.

        \paragraph{Kommunikationsaufbau} Zu Beginn der Laufzeit des Client, wird wie bereits bei der 
        Serverkomponente ein Socket definiert. Über die Angabe des gewünschten Hostnamens und des Ports
        kann die Verbindung mit dem Socket auf dem Server aufgebaut werden.

        Nach dem Aufbau der Verbindung beginnt die Realisierung des Informationsaustauschs zwischen
        Client und Server.
        Alle weiteren Operationen werden in der nachfolgenden While-Schleife abgearbeitet.

        \paragraph{Verarbeiten der Nachrichten} Die vorliegende While-Schleife wird solange ausgeführt,
        bis entweder der Client oder der Server die Verbindung abbrechen. Dies gewährleistet die
        Möglichkeit der erneuten Authentifizierung.

        In der Schleife werden Nachrichten auf Basis der Informationen, die von der SIM-Karte geliefert werden
        an den Server gesendet. Ebenso werden von selbigem Informationen empfangen. Je nach aktuellem Zustand
        bzw. Position im Authentifizierungsvorgang, werden unterschiedliche Operationen ausgeführt.
        Welche Operationen ausgeführt werden müssen, wird anhand des Statuscodes erkannt. Dieser kann
        jeweils aus den Nachrichten des Servers extrahiert werden. Auf entsprechende Werte wird immer
        über die Länge bzw. Start- und Endstelle zugegriffen. So extrahiert der Client beispielsweise
        den Statuscode aus den Stellen 0-1 der Antwortnachricht.

        \paragraph{Beenden der Verbindung} Die Verbindung wird dann beendet, wenn entweder Client
        oder Server dies initiieren. 
	
    \subsubsection{Kommunikationsablauf zwischen Client und Server}
    Dieser Ablauf wird über die durch Sockets initiierte Verbindung realisiert. Aufgrund
    der Tatsache, dass der Vorgang je nach Ausgangszustand den Aufruf verschiedener Operationen
    verlangt, wird ein Statuscode in jede Nachricht integriert. Anhand dieses Statuscodes
    können dann beide Kommunikationsteilnehmer jeweils erforderliche Operationen ausführen.

    Der Client bleibt so lange in der Authentisieurngsschleife, bis diese erfolgreich abgeschlossen werden.
    Gleiches gilt für den Server. Beide reagieren je nach Statuscode (extrahiert aus den ersten Stellen)
    mit nachfolgend erläuterten Aktionen.

    \paragraph{Statuscode 0:} Bisher wurde nur die Verbindung zwischen beiden Teilnehmern erfolgreich aufgebaut,
    aber keine Informationen ausgetauscht. Der Client liest die IMSI der SIM-Karte aus und schickt sie mit
    vorangestelltem \textit{Statuscode 1} an den Server.

    \paragraph{Statuscode 1:} Der Server kann aufgrund des Statuscodes die Beschaffenheit
    der Clientnachricht assoziieren und daraufhin die IMSI extrahieren. Der Client führt
    nach Empfang der Serverantwort, die einen Wert für sowohl AUTN als auch RAND enthält, den
    ersten Authentifizierungsversuch durch. Glückt dieser, generiert die SIM-Karte bereits zu diesem
    Zeitpunkt schon Werte für die Vektoren Ck, Ik und RES und der Authentifizierungsvorgang ist
    erfolgreich beendet. Sollte dies nicht der Fall sein, antwortet die SIM-Karte mit einem Wert für
    AUTS, da die Resynchronisation der SQN noch aussteht. Der Statuscode des Clients wird auf 2 gesetzt
    und eine weitere Iteration des Vorgangs aufgerufen.

    \paragraph{Statuscode 2:} In dieser Iteration sendet der Client neben der IMSI auch die AUTS an den
    Server. Zusätzlich wird dem Server der Statuscode 3 übermittelt.

    \paragraph{Statuscode 3:} Erhält der Server diesen Statusocde führt er die Resynchronisation der SQN
    aus. Dies erfolgt anhand der vom Client erhaltenen AUTS. Hat der Server dies abgeschlossen, übergibt er
    den neu errechneten Wert an den Client zurück. Dieser nimmt den Wert an und führt die Authentifizierung
    erneut aus. Glückt dieser Vorgang antwortet die SIM-Karte schließlich mit den gewünschten Vektoren
    Ck, Ik und RES. Die Authentifizierung ist hier erfolgreich abgeschlossen.

\clearpage
